% =============================================================================
% Derivation of the q(tau) Variational Update
% Self-Contained Reference for Supervised Bayesian MF Implementation
%
% Author:  Andrew Walther
% Date:    March 2026
%
% PURPOSE:
%   Derives the tau (noise precision) update.  Unlike L, F, and beta — all
%   EBNM problems — tau is updated via a CLOSED-FORM MLE that operates on the
%   full matrices simultaneously (no per-k loop).  A critical variance
%   correction term prevents systematic underestimation of noise.
%
% COMPANION CODE: code/update_tau.R
% PARENT DOCUMENT: derivations/MF_UpdateDerivations/MF_Derivations_UpdateAlgo_REVISED.tex
% =============================================================================

\documentclass[11pt]{article}
\usepackage[margin=1in]{geometry}
\usepackage{amsmath, amssymb, amsthm}
\usepackage{bm}
\usepackage{xcolor}
\usepackage{hyperref}
\usepackage{booktabs}
\usepackage{array}
\usepackage{cancel}
\usepackage{tcolorbox}
\tcbuselibrary{skins, breakable}

% ---- Custom commands --------------------------------------------------------
\newcommand{\lbar}[2]{\bar{l}_{#1#2}}
\newcommand{\fbar}[2]{\bar{f}_{#1#2}}
\newcommand{\lsec}[2]{\overline{l^2_{#1#2}}}
\newcommand{\fsec}[2]{\overline{f^2_{#1#2}}}
\newcommand{\revised}[1]{{\color{red}\textbf{[REVISED]:} #1}}

% ---- Coloured boxes ---------------------------------------------------------
\newtcolorbox{correctionbox}{
  colback=red!5!white, colframe=red!50!black,
  title=Correction Note, fonttitle=\bfseries\small, breakable
}
\newtcolorbox{derivbox}{
  colback=gray!6!white, colframe=gray!50!black,
  title=Derivation Steps, fonttitle=\bfseries\small, breakable
}
\newtcolorbox{verifybox}{
  colback=green!5!white, colframe=green!50!black,
  title=Verification Check, fonttitle=\bfseries\small, breakable
}
\newtcolorbox{codebox}{
  colback=yellow!5!white, colframe=orange!50!black,
  title=Code Mapping (\texttt{update\_tau.R}), fonttitle=\bfseries\small, breakable
}
\newtcolorbox{notebox}{
  colback=blue!5!white, colframe=blue!40!black,
  title=Key Property, fonttitle=\bfseries\small, breakable
}

% ---- Document ---------------------------------------------------------------
\title{\textbf{Derivation of the $q(\tau)$ Variational Update}\\[4pt]
       \large Supervised Bayesian Matrix Factorization\\[4pt]
       \normalsize Self-Contained Reference for \texttt{code/update\_tau.R}}
\author{Andrew Walther}
\date{March 2026}

\begin{document}
\maketitle

\begin{abstract}
This document derives the update for the feature-specific noise precision
parameters $\tau_j$ in the Supervised Bayesian Matrix Factorization model.
Unlike the updates for $L$, $F$, and $\beta$ — which are all Empirical Bayes
Normal Means (EBNM) problems solved via the \texttt{ebnm} R package — the
$\tau$ update is a \emph{closed-form maximum likelihood estimate} (MLE).

The key mathematical ingredient is the \emph{expected squared residual}
$\bar{R}^2_{ij}$, which includes a variance correction term that accounts
for posterior uncertainty in $L$ and $F$.  Ignoring this correction leads to
systematic overestimation of $\tau_j$ (underestimation of noise).  The $\tau$
update also produces the genomics ELBO proxy, used as a convergence monitor.

This is the fourth and final modular derivation document
(q$\beta \to$ qL $\to$ qF $\to$ q$\tau$ (this)).
\end{abstract}

\tableofcontents
\newpage

% =============================================================================
\section{Setup and Notation}\label{sec:setup}
% =============================================================================

\subsection{Model and Role of $\tau$}

\begin{equation}
  Y_{ij} = \sum_{k=1}^K l_{ik}f_{jk} + \varepsilon_{ij},
  \qquad \varepsilon_{ij} \sim \mathcal{N}(0,\,\tau_j^{-1})
  \label{eq:genomics}
\end{equation}

Each feature $j$ has its own noise precision $\tau_j > 0$.  A flat
(improper) prior is placed on $\tau_j$, so the update is a pure MLE.
$\tau_j$ does \emph{not} appear in the Cox survival term.

\subsection{Notation}

\begin{center}
\begin{tabular}{lll}
\toprule
\textbf{Symbol} & \textbf{Definition} & \textbf{R variable} \\
\midrule
$\lbar{i}{k}$ & $\mathbb{E}_q[l_{ik}]$ & \texttt{EL[i,k]} \\
$\lsec{i}{k}$ & $\mathbb{E}_q[l_{ik}^2]$ & \texttt{EL2[i,k]} \\
$\fbar{j}{k}$ & $\mathbb{E}_q[f_{jk}]$ & \texttt{EF[j,k]} \\
$\fsec{j}{k}$ & $\mathbb{E}_q[f_{jk}^2]$ & \texttt{EF2[j,k]} \\
$\hat{Y}_{ij}$ & Predicted mean $\sum_k\lbar{i}{k}\fbar{j}{k}$ & \texttt{EL \%*\% t(EF)} \\
$\bar{R}^2_{ij}$ & Expected squared residual & \texttt{R2\_bar[i,j]} \\
$V_{ij}$ & Variance correction term & \texttt{Var\_Term[i,j]} \\
\bottomrule
\end{tabular}
\end{center}

% =============================================================================
\section{ELBO Terms Involving $\tau_j$}\label{sec:elbo}
% =============================================================================

With a flat prior on $\tau_j$, the ELBO contribution from feature $j$ is:
\begin{equation}
  \mathcal{F}_j(\tau_j) =
  \mathbb{E}_q\!\left[\log P(Y_{\cdot j} \mid L, F, \tau_j)\right]
  = \mathbb{E}_q\!\left[\sum_{i=1}^n \left(\frac{1}{2}\log\tau_j
    - \frac{\tau_j}{2}(Y_{ij} - \eta_{ij})^2\right)\right]
  \label{eq:elbo-j}
\end{equation}

where $\eta_{ij} = \sum_k l_{ik}f_{jk}$.  Taking the expectation:
\begin{equation}
  \mathcal{F}_j(\tau_j) = \frac{n}{2}\log\tau_j
  - \frac{\tau_j}{2}\sum_i\mathbb{E}_q\!\left[(Y_{ij} - \eta_{ij})^2\right]
  \label{eq:elbo-j2}
\end{equation}

The expected squared residual $\bar{R}^2_{ij} = \mathbb{E}_q[(Y_{ij}-\eta_{ij})^2]$
is derived in Section~\ref{sec:r2bar}.  Differentiating \eqref{eq:elbo-j2}
with respect to $\tau_j$ and setting to zero:
\begin{equation}
  \frac{d\mathcal{F}_j}{d\tau_j} = \frac{n}{2\tau_j}
  - \frac{1}{2}\sum_i\bar{R}^2_{ij} = 0
  \quad\Rightarrow\quad
  \boxed{\hat{\tau}_j = \frac{n}{\displaystyle\sum_{i=1}^n \bar{R}^2_{ij}}}
  \label{eq:tau-mle}
\end{equation}

% =============================================================================
\section{Expected Squared Residual $\bar{R}^2_{ij}$}\label{sec:r2bar}
% =============================================================================

\subsection{Naive Residual (Insufficient)}

If we naively plugged in posterior means:
\begin{equation}
  (Y_{ij} - \hat{Y}_{ij})^2
  = \left(Y_{ij} - \sum_k\lbar{i}{k}\fbar{j}{k}\right)^2
  \label{eq:naive}
\end{equation}
this ignores posterior uncertainty in $L$ and $F$.

\subsection{Full Expectation}

The correct expected squared residual is:
\begin{align}
  \bar{R}^2_{ij}
  &= \mathbb{E}_q\!\left[\left(Y_{ij} - \sum_k l_{ik}f_{jk}\right)^2\right]
  \nonumber\\
  &= \left(Y_{ij} - \sum_k\lbar{i}{k}\fbar{j}{k}\right)^2
     + \mathbb{E}_q\!\left[\left(\sum_k(l_{ik}f_{jk} - \lbar{i}{k}\fbar{j}{k})\right)^2\right]
  \label{eq:r2bar-split}
\end{align}

The second term (the variance correction) expands via independence of rows of $L$ and $F$ under $q$:

\begin{derivbox}
Assuming $q$ factorises over $(i,k)$ pairs (mean-field):
\begin{align*}
  \mathbb{E}_q\!\left[\left(\sum_k l_{ik}f_{jk}\right)^2\right]
  &= \sum_k\mathbb{E}_q[l_{ik}^2]\,\mathbb{E}_q[f_{jk}^2]
  \quad\text{(cross terms vanish by independence)}\\
  &= \sum_k\lsec{i}{k}\,\fsec{j}{k}
\end{align*}

Meanwhile:
\begin{align*}
  \left(\mathbb{E}_q\!\left[\sum_k l_{ik}f_{jk}\right]\right)^2
  &= \left(\sum_k\lbar{i}{k}\fbar{j}{k}\right)^2
  = \sum_k\lbar{i}{k}^2\fbar{j}{k}^2
    + 2\sum_{k<k'}\lbar{i}{k}\fbar{j}{k}\lbar{i}{k'}\fbar{j}{k'}
\end{align*}

Therefore the variance of the prediction is:
\begin{align}
  \mathrm{Var}_q(\eta_{ij})
  &= \mathbb{E}_q[\eta_{ij}^2] - (\mathbb{E}_q[\eta_{ij}])^2
  \nonumber\\
  &= \sum_k\lsec{i}{k}\,\fsec{j}{k}
     - \left(\sum_k\lbar{i}{k}\fbar{j}{k}\right)^2
  \label{eq:var-pred}
\end{align}

But this is still not the right expansion.  Using the identity
$\mathbb{E}[(X-c)^2] = \mathrm{Var}(X) + (\mathbb{E}[X]-c)^2$:

\begin{align}
  \bar{R}^2_{ij}
  &= (Y_{ij} - \hat{Y}_{ij})^2 + \mathrm{Var}_q(\eta_{ij})
  \nonumber\\
  &= (Y_{ij} - \hat{Y}_{ij})^2
     + \underbrace{\sum_k\lsec{i}{k}\,\fsec{j}{k}
       - \sum_k\lbar{i}{k}^2\fbar{j}{k}^2}_{V_{ij}}
  \label{eq:r2bar-final}
\end{align}
\end{derivbox}

\subsection{Variance Correction Term $V_{ij}$}

\begin{equation}
  \boxed{V_{ij} = \sum_k \bigl(\lsec{i}{k}\,\fsec{j}{k}
                              - \lbar{i}{k}^2\,\fbar{j}{k}^2\bigr)
  = (EL2\,EF2^\top)_{ij} - (EL^2\,EF^{2\top})_{ij}}
  \label{eq:var-term}
\end{equation}

\textbf{Non-negativity:} Each term $\lsec{i}{k}\,\fsec{j}{k} - \lbar{i}{k}^2\fbar{j}{k}^2 \geq 0$
because $\mathbb{E}[X^2]\mathbb{E}[Y^2] \geq (\mathbb{E}[X]\mathbb{E}[Y])^2$ by
independence and Jensen's inequality.  Therefore $V_{ij} \geq 0$ and
$\bar{R}^2_{ij} \geq (Y_{ij}-\hat{Y}_{ij})^2$.

\textbf{Consequence:} The variance correction \emph{inflates} the expected
squared residual, which \emph{reduces} $\hat{\tau}_j$ relative to the naive
estimate.  Omitting $V_{ij}$ causes systematic overestimation of precision
(underestimation of noise).

\begin{codebox}
\begin{verbatim}
compute_var_term <- function(EL, EL2, EF, EF2) {
  (EL2 %*% t(EF2)) - (EL^2 %*% t(EF^2))
}
\end{verbatim}

\textbf{Matrix form explanation:}
\begin{itemize}
  \item \texttt{EL2 \%*\% t(EF2)}: entry $(i,j)$ is $\sum_k\lsec{i}{k}\fsec{j}{k}$
  \item \texttt{EL\^{}2 \%*\% t(EF\^{}2)}: entry $(i,j)$ is $\sum_k\lbar{i}{k}^2\fbar{j}{k}^2$
\end{itemize}
\end{codebox}

\begin{correctionbox}
\textbf{Corrections R7 and R8} (from REVISED.tex):
\begin{itemize}
  \item R7: Sign error — the variance correction term enters with a
    \emph{positive} sign in $\bar{R}^2_{ij}$, not negative.
    The original document had an erroneous minus sign.
  \item R8: Subscript — the loading index should be $l_{ik}$ (sample $i$,
    factor $k$), not $l_{jk}$ which refers to a feature.
\end{itemize}
Both are fixed in the REVISED document and reflected in \texttt{code/update\_tau.R}.
\end{correctionbox}

% =============================================================================
\section{The Full Tau Update}\label{sec:update}
% =============================================================================

Combining \eqref{eq:tau-mle} and \eqref{eq:r2bar-final}:

\begin{align}
  \bar{R}^2_{ij} &= (Y_{ij} - \hat{Y}_{ij})^2 + V_{ij} \label{eq:r2bar-combined} \\[6pt]
  \hat{\tau}_j &= \frac{n}{\sum_i \bar{R}^2_{ij}} \label{eq:tau-hat}
\end{align}

In matrix notation:
\begin{equation}
  R2\_bar = (Y - EL\,EF^\top)^2 + V
  \qquad\text{(element-wise squared + correction)}
  \label{eq:r2bar-matrix}
\end{equation}
\begin{equation}
  \hat{\bm{\tau}} = n\,/\,\mathrm{colSums}(R2\_bar)
  \qquad\text{(element-wise division)}
  \label{eq:tau-vector}
\end{equation}

\begin{notebox}
\textbf{Why this is NOT an EBNM problem.}

The updates for $L$, $F$, and $\beta$ all involve placing a prior $g(\cdot)$
on the parameter and running an EBNM solver.  The $\tau$ update is different:
\begin{itemize}
  \item $\tau_j$ has a flat (improper uniform) prior — no shrinkage.
  \item The optimisation is unconstrained (no prior KL term).
  \item The MLE has a closed form: $\hat{\tau}_j = n/\sum_i\bar{R}^2_{ij}$.
  \item There is no per-$k$ loop: all $K$ factors contribute simultaneously
    through $\hat{Y}_{ij} = \sum_k\lbar{i}{k}\fbar{j}{k}$ and $V_{ij}$.
\end{itemize}
\end{notebox}

\begin{codebox}
\begin{verbatim}
update_tau <- function(Y, EL, EL2, EF, EF2, tau_floor = 1e-8) {
    n        <- nrow(Y)
    Var_Term <- compute_var_term(EL, EL2, EF, EF2)          # n x p
    R2_bar   <- (Y - EL %*% t(EF))^2 + Var_Term             # n x p
    col_sums <- colSums(R2_bar)                              # p-vector
    Tau      <- n / pmax(col_sums, n * tau_floor)            # p-vector
    ...
}
\end{verbatim}
The \texttt{pmax} floor prevents $\hat{\tau}_j = \infty$ when the residual
sum is near zero (perfect reconstruction).
\end{codebox}

% =============================================================================
\section{ELBO Proxy as Convergence Monitor}\label{sec:elbo-proxy}
% =============================================================================

Substituting the MLE $\hat{\tau}_j$ back into \eqref{eq:elbo-j2}:

\begin{equation}
  \mathcal{F}_j(\hat{\tau}_j)
  = \frac{n}{2}\log\hat{\tau}_j
  - \frac{\hat{\tau}_j}{2}\sum_i\bar{R}^2_{ij}
  = \frac{n}{2}\log\hat{\tau}_j - \frac{n}{2}
  \label{eq:elbo-at-mle}
\end{equation}

The genomics \emph{ELBO proxy} (summed over all features) is:
\begin{equation}
  \boxed{\mathcal{E}_{\text{proxy}} = \sum_j\left[\frac{n}{2}\log\hat{\tau}_j
  - \frac{\hat{\tau}_j}{2}\sum_i\bar{R}^2_{ij}\right]}
  \label{eq:elbo-proxy}
\end{equation}

This is the genomics contribution $\mathbb{E}_q[\log P(Y \mid L, F, \tau)]$
evaluated at the current posterior means and the updated $\tau$.  It is
monotonically increasing (or at worst non-decreasing) across CAVI iterations
if the algorithm is working correctly.

\begin{codebox}
\begin{verbatim}
elbo_proxy <- sum(n / 2 * log(Tau) - Tau / 2 * col_sums)
\end{verbatim}
\end{codebox}

% =============================================================================
\section{Verification Checks}\label{sec:verification}
% =============================================================================

\begin{verifybox}
\textbf{V1: Non-negativity.}
$V_{ij} \geq 0$ (element-wise), therefore $\bar{R}^2_{ij} \geq (Y_{ij} -
\hat{Y}_{ij})^2 \geq 0$.  If $V_{ij}$ contains negative entries, there is a
sign error in the implementation.

\textbf{V2: Perfect reconstruction.}
When $Y = EL\,EF^\top$ exactly and $EL2 = EL^2$, $EF2 = EF^2$:
$\bar{R}^2_{ij} = 0$, so $\hat{\tau}_j$ hits the ceiling $1/\tau_{\text{floor}}$.

\textbf{V3: Correction direction.}
Ignoring $V$ (using $EL2 = EL^2$, $EF2 = EF^2$) gives $V_{ij} = 0$ and
hence a \emph{larger} $\hat{\tau}_j$ (overestimates precision).  The
corrected estimate is always $\leq$ the biased estimate.

\textbf{V4: ELBO increases.}
Over a sequence of CAVI iterations (L update $\to$ F update $\to$ $\tau$
update), the ELBO proxy should be non-decreasing.  A decrease signals a bug
in one of the update equations.

\textbf{V5: V2.R consistency.}
The implementation should match V2.R lines 374--376 exactly:
\begin{center}
\texttt{Var\_Term = EL2 \%*\% t(EF2) - EL\^{}2 \%*\% t(EF\^{}2)}\\
\texttt{Tau = n / pmax(colSums(R2\_bar), n * tau\_floor)}
\end{center}
\end{verifybox}

% =============================================================================
\section{Code Mapping}\label{sec:code}
% =============================================================================

\begin{center}
\begin{tabular}{lll}
\toprule
Equation & Code & Type \\
\midrule
$V_{ij} = \sum_k(\lsec{i}{k}\fsec{j}{k} - \lbar{i}{k}^2\fbar{j}{k}^2)$ & \texttt{compute\_var\_term()} & $n\times p$ \\
$\bar{R}^2_{ij} = (Y-\hat{Y})^2_{ij} + V_{ij}$ & \texttt{R2\_bar} & $n\times p$ \\
$\hat{\tau}_j = n/\sum_i\bar{R}^2_{ij}$ & \texttt{n / pmax(col\_sums, ...)} & $p$-vector \\
$\mathcal{E}_{\text{proxy}}$ & \texttt{elbo\_proxy} & scalar \\
\bottomrule
\end{tabular}
\end{center}

% =============================================================================
\section{Algorithm Summary}\label{sec:algorithm}
% =============================================================================

\begin{enumerate}
  \item Compute variance correction (matrix operation, all $K$ factors at once):
    \[V = EL2\,EF2^\top - EL^2\,(EF^2)^\top \quad[n\times p]\]
  \item Compute expected squared residuals:
    \[\bar{R}^2 = (Y - EL\,EF^\top)^{\odot 2} + V \quad[n\times p]\]
  \item Compute column sums and update precision:
    \[\hat{\tau}_j = \max\!\left(\frac{n}{\displaystyle\sum_i\bar{R}^2_{ij}},\;
    \frac{1}{\tau_{\text{floor}}}\right)^{-1} \quad[p\text{-vector}]\]
  \item Compute ELBO proxy for convergence monitoring:
    \[\mathcal{E}_{\text{proxy}} = \sum_j\left[\frac{n}{2}\log\hat{\tau}_j
      - \frac{\hat{\tau}_j}{2}\sum_i\bar{R}^2_{ij}\right]\]
  \item Return $\hat{\bm{\tau}}$, $\bar{R}^2$, $V$, $\mathcal{E}_{\text{proxy}}$.
\end{enumerate}

% =============================================================================
\appendix
\section{Notation Index}
% =============================================================================

\begin{tabular}{lll}
\toprule
Symbol & Meaning & Dimensions \\
\midrule
$Y$ & Genomics data & $n \times p$ \\
$\hat{Y}_{ij}$ & $\sum_k\lbar{i}{k}\fbar{j}{k}$ & $n \times p$ \\
$V_{ij}$ & Variance correction term ($\geq 0$) & $n \times p$ \\
$\bar{R}^2_{ij}$ & Expected squared residual & $n \times p$ \\
$\hat{\tau}_j$ & Updated noise precision & $p$-vector \\
$\mathcal{E}_{\text{proxy}}$ & Genomics ELBO proxy & scalar \\
$\tau_{\text{floor}}$ & Precision floor (default $10^{-8}$) & scalar \\
\bottomrule
\end{tabular}

\end{document}
